\section{A coupled fermion current in a declared hybrid gauge action}
\label{app:fermion-coupled-action}

The fermionic current can be coupled to electric drift and magnetic
plaquette forces on the same prepared carrier. The construction below
declares a regulated Hamiltonian and a finite splitting word. Its gauge
fields are classical, while its matter state belongs to finite fermion
Fock space. This hybrid interpretation differs from quantizing the gauge
field in the preceding one-link construction.

Take the direct-product gauge cover
$SU(3)\times SU(2)\times U(1)$, with the five registered representations
and integer hypercharges $q_s=6Y_s$ in each of three declared families.
The compact abelian link is $U_e=e^{i\theta_e}$ with momentum $E_e$.
For fixed $M,\lambda>0$, choose
\begin{align}
 T_\lambda(E)&=\frac{2E}{\lambda}
       \arctan\frac{\lambda E}{2M}
       -\frac{2M}{\lambda^2}
       \log\left(1+\frac{\lambda^2E^2}{4M^2}\right),\\
 W_f&=\prod_{e\in\partial f}U_e^{\epsilon_{fe}},\qquad
 V_f=\beta_f(1-\operatorname{Re}W_f).
\end{align}
Here $\epsilon_{fe}$ is the oriented face boundary. The reduced Hamiltonian
is
\[
 \mathcal H_{\mathrm{red}}=
 \sum_e T_\lambda(E_e)+\sum_f V_f+
 \sum_e\langle\Psi,\widehat H_{\mathrm{CAR},e}(U_e)\Psi\rangle,
\]
where the hopping operators and Spin matrices are those defined above.
In temporal gauge, its first-order action contains
$\sum_e E_e\dot\theta_e+i\langle\Psi,\dot\Psi\rangle-
\mathcal H_{\mathrm{red}}$. The link-phase variation of the same matter
Hamiltonian defines its current.

This is a specified kinetic law. In particular,
\[
 T_\lambda'(E)=\frac{2}{\lambda}\arctan\frac{\lambda E}{2M},
 \qquad
 T_\lambda''(E)=\frac{1}{M(1+(\lambda E/(2M))^2)}>0.
\]
Integrating $x-\arctan x=\int_0^x t^2/(1+t^2)\,dt$ gives
\[
 0\leq\frac{E^2}{2M}-T_\lambda(E)
       \leq\frac{\lambda^2E^4}{48M^3}.
\]
The quadratic electric energy is therefore a controlled small-field
comparison of functions. It is not the exact kinetic energy of this model,
and the displayed bound is not a trajectory-error estimate. The parameter
$\lambda$ remains fixed when describing this action; varying it changes
the law. The separate charge-one abelian Wilson term is specified on the
direct-product cover. No descent of that term to the global central
quotient is assumed.

\begin{proposition}[Invariant expectation-level matter sector]
Extend the reduced Hamiltonian by classical nonabelian links and momenta,
quadratic nonabelian electric energies, Wilson magnetic energies,
the covariant Higgs kinetic and spatial terms, and a potential with vanishing
first derivative at zero Higgs field. Couple the finite CAR matter through
its registered representation matrices and all twenty-seven complex Yukawa
coefficients of three families. Interpret every matter force as its Fock-state
expectation in the canonical classical-field/quantum-matter action.

Set the nonabelian links to identity, their electric momenta to zero,
and the Higgs field and conjugate momentum to zero. In each species of one
family take normal covariance
$P_s=I_{d_s}\otimes|\psi_s\rangle\langle\psi_s|$,
with no covariance between distinct species or families. Let the other two
families be Fock vacuum, and take zero anomalous covariance
$\langle c_i c_j\rangle=0$. These conditions define an invariant hybrid
subsystem with Hamiltonian $\mathcal H_{\mathrm{red}}$.
\end{proposition}

\begin{proof}
Every nonabelian matter charge or link-current variation inserts a
traceless internal generator $T_s$. The remaining Spin and spatial factor
is scalar in that internal index, so the expectation factors through
$\operatorname{tr}T_s=0$. This checks the omitted nonabelian matter forces,
including each color and weak generator. The nonabelian electric velocities
vanish at zero momenta, and their magnetic first variations vanish at
identity plaquette holonomy. Scalar kinetic, spatial and gauge-current
variations vanish at zero Higgs field and momentum.

In all-left-Weyl notation, the Yukawa terms have the gauge-neutral forms
$QHu^c$, $Q\overline H d^c$ and $L\overline H e^c$, with the required
Spin, color and weak contractions and their adjoints. Their Higgs first
variation is a linear combination of anomalous expectations
$\langle cc\rangle$ and $\langle c^\dagger c^\dagger\rangle$.
Both vanish by fermion-number selection, for every value of every Yukawa
coefficient. This checks the Higgs force rather than only the value of the
Yukawa energy at $H=0$.

The surviving hopping Hamiltonian is number conserving, family diagonal
and scalar on each internal multiplicity space. It preserves the stated
covariance form, the vanishing anomalous covariance and the two vacuum
families. The field equations hence preserve every imposed condition.
No stability of the zero-Higgs equilibrium is needed for this invariance
statement.
\end{proof}

The expectation hypothesis is essential. It does not identify independent
Grassmann components with one commuting orbital, and it does not annihilate
the nonabelian Gauss operators. For example, in the initial uniform
preparation $p=1/64$, the weak generator $\sigma_z/2$ on the occupied quark
and lepton doublets has local charge variance
\[
 \operatorname{Var}(\widehat Q_T(v))=
 \sum_s\operatorname{tr}(T_s^2)p(1-p)
 =\frac{63}{2048}>0.
\]
The absent nonabelian mean source therefore cannot certify a full
gauge-constrained quantum state with sharp zero nonabelian electric data.

\begin{proposition}[Coupled transport and conserved readout]
On a finite oriented lattice with closed face boundaries, compose the
following factors of the declared hybrid action:
\begin{enumerate}
\item a fermionic implicit-midpoint hopping factor, with its variational
      current and the same classical electric kick;
\item an electric factor of duration $\lambda$,
\[
 U_e'=
 \frac{1+i\lambda E_e/(2M)}{1-i\lambda E_e/(2M)}U_e,
 \qquad E_e'=E_e;
\]
\item a magnetic factor of duration $h$ on face $f$,
\[
 E_e'=E_e-h\beta_f\epsilon_{fe}\operatorname{Im}W_f
 \quad(e\in\partial f),\qquad U_e'=U_e.
\]
\end{enumerate}
For Gaussian-rational initial data and rational parameters these maps are
exact over Gaussian rationals. Each preserves the expectation Gauss residual
at every vertex. Their recorded current and charge transfers satisfy exact
discrete continuity under finite composition, and the construction is
covariant under residual time-independent vertex gauge transformations.
\end{proposition}

\begin{proof}
The midpoint matter identity proved above gives opposite endpoint charge
transfers and precisely the electric kick that cancels them in Gauss.
The midpoint equations are applied to the orbital, link-angle and momentum
variables of that same factor action; its angle is constant during the
factor because its Hamiltonian has no electric-momentum dependence.

The electric factor is the exact flow of $T_\lambda$ for the stated
duration: exponentiating
$i\lambda T_\lambda'(E)=2i\arctan(\lambda E/(2M))$
gives the rational phase. It leaves charge and electric divergence
unchanged. For the magnetic factor,
$\partial_{\theta_e}V_f=\beta_f\epsilon_{fe}\operatorname{Im}W_f$.
Its momentum kick is the oriented face boundary, whose vertex divergence
vanishes. These facts prove Gauss preservation for every factor and hence
for their composition. The unit-modulus rational phase, Hermitian Cayley
matter transform and rational Wilson product also prove the asserted
arithmetic closure. Abelian gauge transformations conjugate the matter
updates, leave $E$ and $W_f$ invariant and commute with the electric drift.
\end{proof}

This factor word is a declared discrete transport. It is not the exact
exponential of the summed Hamiltonian; exact Gauss conservation does not
imply exact total-energy conservation or a continuous-trajectory error bound.

\paragraph{Finite execution and returned current.}
Use the same sixty-four prepared addresses, 144 oriented links and 108
plaquettes, with $\lambda=h=1/2$ and $M=\beta=\kappa=1$. A full matter
sweep is followed by all electric drifts and all magnetic kicks. One
subsequent matter factor reads the evolved link and momentum on edge~14.
That final factor is a readout in a partial second word, not a second
completed full-Hamiltonian time step.

Across baseline, local intervention, gauge copy and a disabled-drift
control, the retained prefixes contain 2,444 parent events.
Each of the four continuations has 255 events: 144 electric factors, 108
magnetic factors, one returning matter factor and two public readouts.
The control leaves the link unchanged instead of applying electric drift;
it is a counterfactual operation word, not another execution of the full
declared action. In the intervention and gauge copy the return current lies
in $[-0.005355160745,-0.005355160744]$, and the final edge momentum lies in
$[0.002253305815,0.002253305816]$. The difference from the disabled-drift
return current obeys the exact outward enclosure
\[
 0<\frac{242747271}{2^{40}}
 \leq J_{\mathrm{return}}-J_{\mathrm{disabled}}
 \leq\frac{30343409}{2^{37}}.
\]
Thus the electric record changes the link that the later fermionic current
actually consumes. Gauge-related executions agree on the invariant readouts.

The package \texttt{code/sm\_fermion\_current/} contains the producer
\texttt{coupled.py} and independent verifier \texttt{verify\_coupled.py}.
They reconstruct the closed face boundaries, exact factor updates, current,
Gauss residuals and consumed writer/value versions. The observer-like
software patches retain local orbital state, incident ports, public records
and electric feedback; declared parent labels cannot substitute for actual
consumption. Local magnetic energies are exact rationals, with their large
sum bounded by outward dyadic intervals.

The action, regulator, family population, prepared addresses, routing and
time labels are supplied. The invariant reduction does not execute an
excited Higgs or nonabelian sector, select a chiral continuum regulator,
identify a laboratory current or physical clock, or establish compatibility
with a sourced curved geometry. Those statements require their own
common-action and physical realization maps.
